\documentclass[11pt]{article}
\usepackage{amsmath,amssymb,amsthm}
\newtheorem{theorem}{Theorem}
\newtheorem{lemma}[theorem]{Lemma}
\title{Coalescing Reciprocal Columns and Graded Flags}
\author{}
\date{November 23, 2026}
\begin{document}
\maketitle

\paragraph{Status and scope.}
This is a complete written audit of one proposed decoder, not a proof of
either global asymptotic target. The statements have not been formalized
in Lean. The exact companion checker verifies the identities and the
nongraded four-column example. The sign-count ingredient is the
unconditional theorem
\texttt{HeadComplexity.strictPolynomialSignPatterns\_card\_le} in
\texttt{Results/WarrenBound.lean}.

\paragraph{The proposed decoder.}
For each slot $h$ one hoped that coalescing columns would supply a feature
$v_h W^{k_h}$, with the orders $k_h$ permuted freely among the slots.
Rows indexed by a token address and a binary digit of the slot would then
have profile dimension only $d=O(\log H)$. Positive residues proportional
to $2^{-Mk^2}$ and addresses $W=2^{2Mk_0}$ would select the term $k=k_0$.
The missing issue is whether the proposed feature space is an actual
limit of the reciprocal columns. The following constraint rules out a
graded first-order version of that step.

\begin{lemma}[A finite-dimensional limit invariant]
Fix vectors $v_1,\ldots,v_J\in\mathbb R^d$ and real poles
$\lambda_h(t)\to0$ as $t\to0$. Consider the vector-valued columns
\[
                 \frac{v_h}{1-\lambda_h(t)W}.
\]
Multiply their span by
$D_t(W)=\prod_h(1-\lambda_h(t)W)$, so it is a subspace $E_t$ of
$\mathbb R^d\otimes\mathbb R[W]_{<J}$.
Suppose $E_t$ has constant dimension $L$ and a Grassmann limit $E_0$.
Then $E_0$ is invariant under the backward shift
\[
                 T p(W)=\frac{p(W)-p(0)}{W}.
\]
The same $E_0$ is the limit of the original rational-function spans:
their common multiplier $D_t$ tends coefficientwise to $1$.
\end{lemma}
\begin{proof}
On the fixed finite-dimensional polynomial space define
\[
          T_t p(W)=\frac{p(W)-p(0)D_t(W)}{W}.
\]
The quotient is polynomial of degree less than $J$. Its matrix tends
to that of $T$. For the $h$th numerator
$p_h=v_hD_t/(1-\lambda_hW)$, direct substitution gives
$T_t p_h=\lambda_h p_h$. Thus $T_tE_t\subseteq E_t$.
Choose bases for $E_t$ converging to a basis for $E_0$, which is a
standard local coordinate description of a Grassmann limit. Passing
to the limit gives $TE_0\subseteq E_0$.
\end{proof}

\paragraph{Reciprocal powers.}
The analogous statement for $v_h/(1-\lambda_hW)^q$, with a fixed positive
integer $q$, uses the weighted shift
\[
       T_q(W^k)=\frac{k}{k+q-1}W^{k-1}\quad(k\geq1),
       \qquad T_q(1)=0.
\]
Indeed, the coefficient multiplier
$J_q(W^k)=\binom{k+q-1}{k}W^k$ carries a simple reciprocal to its
$q$th power and conjugates $T$ to $T_q$. For integral $q$, $J_q$ is
a polynomial in the Euler derivative $W\,d/dW$, so passage to a
finite polynomial limit commutes with $J_q$. In particular the
first-order profile terms from a reciprocal perturbation use $q=2$.

\begin{theorem}[Graded limits have nested profile spaces]
Suppose an $L$-dimensional invariant limit is graded:
\[
                  E_0=\bigoplus_{k\geq0} S_k W^k,
                  \qquad S_k\subseteq\mathbb R^d.
\]
Then $S_{k+1}\subseteq S_k$. There are a basis $u_1,\ldots,u_d$ of
$\mathbb R^d$ and integers
$\ell_1\geq\cdots\geq\ell_d\geq0$, with $\sum_j\ell_j=L$, such that
every element of $E_0$ has the form
\[
                  p(W)=\sum_{j=1}^d u_j P_j(W),
                  \qquad \deg P_j<\ell_j.
                  \tag{1}
\]
The same conclusions hold with $T_q$ in place of $T$.
\end{theorem}
\begin{proof}
Applying the shift to $vW^{k+1}$ gives a nonzero scalar multiple of
$vW^k$, proving the inclusion. Choose a basis adapted to the finite
nested chain of spaces $S_k$. Let $\ell_j$ be the number of levels
containing its $j$th basis vector. This gives (1), and the dimension
identity gives $\sum_j\ell_j=L$.
\end{proof}

\begin{theorem}[Uniform sign count for all graded flags]
Fix any $M$ rows $(W_i,a_i)\in\mathbb R\times\mathbb R^d$.
For fixed $L,d\geq1$, let $\mathcal F_{L,d}$ contain all scalar row
functions $a_i\cdot p(W_i)$ from every graded, shift-invariant
$L$-dimensional space. Put $p=L+d^2$.
For $M\geq p$, the number of strict sign patterns is at most
\[
           (L+1)^d\left(\frac{8eM}{p}\right)^p.
           \tag{2}
\]
Consequently, the number of rows that can be shattered is
\[
                O\bigl(L+d^2+d\log(L+1)\bigr),
                \tag{3}
\]
with an absolute implied constant, uniformly over the rows.
In particular $d=O(\log L)$ gives only $O(L)$ shattering.
\end{theorem}
\begin{proof}
There are at most $(L+1)^d$ choices of the ordered length tuple
$(\ell_j)$; allowing all tuples with entries between $0$ and $L$
only enlarges the count. For a fixed tuple summing to $L$, allow
arbitrary coordinates for the $d$ vectors $u_j$, using $d^2$ real
parameters, and arbitrary coefficients for the polynomials $P_j$,
using $L$ more. Each row evaluation
\[
       \sum_{j=1}^d\sum_{k=0}^{\ell_j-1}
              (a_i\cdot u_j)c_{jk}W_i^k
\]
is polynomial of total parameter degree at most two. Dropping basis
independence constraints enlarges the family. Warren's strict
sign-pattern estimate, with degree two and $p$ parameters, gives
$(8eM/p)^p$. Summing over tuples proves (2).

For completeness, (2) implies (3) without an extra logarithm in $L$.
Set $b=d\log(L+1)$ and $x=M/p$. If all patterns occur, then
$M\log2\leq b+p\log(8ex)$.
For all sufficiently large absolute $x$, one has
$\log(8ex)\leq x\log2/2$. Thus either $M\leq Cp$, or
$M\leq2b/\log2$. The case $M<p$ is already covered.
\end{proof}

\paragraph{Binary profile powers do not automatically escape.}
For a profile $V_h$ taking only values $\{-1,1\}$, the exact identity
\[
 \frac1{1-\varepsilon W-\delta_h-\eta_hV_h}
 =\frac{1+V_h}{2(1-\varepsilon W-\delta_h-\eta_h)}
  +\frac{1-V_h}{2(1-\varepsilon W-\delta_h+\eta_h)}
 \tag{4}
\]
expresses each column inside the span of two simple reciprocal
columns. After absorbing a scalar, each denominator is
$1-\lambda W$, with
$\lambda=\varepsilon/(1-\delta_h\mp\eta_h)\to0$.
If the profiles $V_h$ span $d$ row profiles, the new profiles
$1\pm V_h$ span at most $d+1$. Thus, if the enlarged $2H$-column
limit is graded, (2)--(3) still give only $O(H)$ shattering when
$d=O(\log H)$. This observation includes all powers of binary
profiles, but it does not establish gradedness.

\paragraph{Gradedness is a real restriction.}
It must not be silently inferred from positivity or coalescence.
Here is an explicit nongraded limit with all poles positive. Take
\[
 v_1=e_1,\quad v_2=e_2,\quad v_3=e_1+e_2,\quad v_4=e_1-e_2,
 \qquad
 (\lambda_1,\lambda_2,\lambda_3,\lambda_4)
 =\left(t,2t,3t,\frac53t+t^2\right).
 \tag{5}
\]
Write $f_h=v_h/(1-\lambda_hW)$,
$g_3=f_3-f_1-f_2$, and $g_4=f_4-f_1+f_2$.
With $u=2e_1+e_2$, four independent column combinations converge to
\[
 \begin{split}
       f_1&\longrightarrow e_1,\qquad
       f_2\longrightarrow e_2,\qquad
       t^{-1}g_3\longrightarrow uW,\\
       t^{-2}(3g_4-g_3)
           &\longrightarrow3(e_1-e_2)W-\frac43uW^2.
 \end{split}
 \tag{6}
\]
This four-dimensional limit is invariant under $T$, as the lemma
requires. It is not graded: the quadratic component $uW^2$ cannot
be separated from the linear component $3(e_1-e_2)W$ using the
other three basis elements. In particular the two coefficient
directions $u$ and $e_1-e_2$ are independent. All denominators in
(5) are positive on every fixed bounded set of $W$ for sufficiently
small positive $t$.

\paragraph{What remains.}
The proposed independently assigned monomial features have not been
constructed. The graded version with binary slot profiles is ruled
out by (2), even if arbitrary graded flags are granted for free.
Nongraded limits, nonbinary profiles of larger span, and correction
families outside the enlarged span in (4) are not covered by that
conclusion. Example (6) shows that investigating nongraded limits
is mathematically substantive. None of these statements supplies
the requested global upper or lower bound.

\paragraph{Arbitrary nongraded invariant spaces can be much larger.}
There is no corresponding $O(L+d^2)$ parameterization of all invariant
spaces. The following explicit family also shows why proving
realizability by the reciprocal columns matters. Suppose $d\geq2$ and
$L\geq d+2$, and put $q=L-d+1$. For arbitrary real coefficients $c_{bk}$,
define
\[
 p_c(W)=e_1W^{q-1}
       +\sum_{b=2}^d\sum_{k=1}^{q-2}c_{bk}e_bW^k,
 \qquad
 E_c=\operatorname{span}\{e_2,\ldots,e_d,p_c,Tp_c,\ldots,T^{q-1}p_c\}.
 \tag{7}
\]
The first components of the last $q$ vectors are the distinct monomials
$W^{q-1},\ldots,1$. Thus the listed vectors are independent,
$\dim E_c=L$, and $E_c$ contains every constant vector.
It is $T$-invariant because the shift sends each listed vector to the
next or to zero.

These spaces are distinct for distinct coefficient arrays. If
$E_c=E_{c'}$, express $p_{c'}$ in the displayed basis of $E_c$.
Its first component forces coefficient one on $p_c$ and zero on the
other shifted vectors. Its zero constant term then forces zero on
$e_2,\ldots,e_d$. Consequently $p_{c'}=p_c$ and $c'=c$.
Hence this family has
$(d-1)(q-2)=(d-1)(L-d-1)$ freely recoverable real parameters.

It also realizes all strict signs on that many fixed rows. Choose
distinct positive $W_1,\ldots,W_{q-2}$ and rows
$(W_j,a=e_b)$ for $2\leq b\leq d$.
The matrix $(W_j^k)_{1\leq j,k\leq q-2}$ is invertible: it is a
Vandermonde matrix times a nonsingular diagonal matrix. Independently
for each $b$, choose the coefficients $c_{bk}$ so that
$e_b\cdot p_c(W_j)$ equals any prescribed value $+1$ or $-1$.

This is an exact shattering statement for the enlarged class of all
invariant spaces. It is not a reciprocal-head construction. In
particular, the fixed-profile maps
\[
 (\lambda_1,\ldots,\lambda_J)
 \longmapsto
 \operatorname{span}\left\{
 v_h\prod_{k\ne h}(1-\lambda_kW):1\leq h\leq J\right\}
 \tag{8}
\]
are highly restricted rational maps into the Grassmannian. Establishing
which members of (7), if any, occur in their coalescing limits is an
additional problem. Neither invariant-space dimension nor the
shattering statement for (7) resolves that problem.
\end{document}
